%% ma: L11-B15-0075
%% lop: 11
%% chuong: 5
%% bai: 15
%% dang: 11.15.5
%% loai: DS
%% muc: [NB, TH, TH, VD]
%% dap_an: "ĐSĐĐ"
%% nguon: "tai-lieu-hoc-tap-mon-toan-lop-11-theo-chuong-trinh-giao-duc-pho-thong-2018.pdf (Lớp Toán Thầy Duy, Phạm Lê Duy, THPT Chu Văn An), Chương 3 Bài 1, Phần 2, Câu 5 (tr.337)"
%% kien_thuc: "Lớp 11 Bài 15 (giới hạn của dãy số)"
%% trang_thai: chinh-thuc
%% ghi_chu: "Đã duyệt 03/11/2026."
\begin{ex}
Cho giới hạn $L=\lim\sqrt{3+\dfrac{an^2-1}{3+n^2}-\dfrac1{2^n}}$. Khi đó
\choiceTF
{\True $L=2$ khi $a=1$.}
{$L=3$ thì có $2$ giá trị nguyên $a$ thỏa mãn.}
{\True $L>3$ khi $a>6$.}
{\True Có $3$ giá trị nguyên của $a$ thuộc $(0;20)$ sao cho $\lim\sqrt{3+\dfrac{an^2-1}{3+n^2}-\dfrac1{2^n}}$ là một số nguyên.}
\loigiai{\sloppy $\dfrac{an^2-1}{3+n^2}\to a$ và $\dfrac1{2^n}\to0$, nên $L=\sqrt{3+a}$ (với $a\ge-3$).
a) $a=1$: $L=2$. Đúng.
b) $L=3\Leftrightarrow3+a=9\Leftrightarrow a=6$: chỉ $1$ giá trị. Sai.
c) $L>3\Leftrightarrow3+a>9\Leftrightarrow a>6$. Đúng.
d) $a\in(0;20)$ nguyên, $L$ nguyên khi $3+a\in\{4;9;16\}$ (vì $3+a\le22$), tức $a\in\{1;6;13\}$: có $3$ giá trị. Đúng.}
\end{ex}
